Over the last twenty years, an extensive literature in \emph{approximate group theory} has been developed establishing \emph{finitary} versions of Gromov's theorem and Trofimov's theorem, highlights of which include \cite{bgt12,shalom-tao,MR2833482,tt.Trof,MR2827010}. See \cite{breuillard2014brief} for a detailed overview, \cite{MR3971253} for a textbook introduction, and \cite{tointon2020brief} for a concise survey.
For groups, this theory culminated in the celebrated work of Breuillard, Green, and Tao \cite{bgt12}, a special case of whose results can be stated\footnote{We state their theorem in a `metric' form that is convenient for our applications, and which is adapted from Tessera and Tointon's structure theorem for \emph{vertex-transitive graphs} of polynomial growth \cite[Theorem 2.3]{tt.Trof}. Indeed, the statement given below is equivalent to the special case of their theorem in which the graph $\Gamma$ is the Cayley graph of $G$, together with the growth bound of \cite[Corollary 11.9]{bgt12}. 
%
} as follows. Given a group $G$ and a finite generating set $S$, we write $\Gr(r)=\Gr_{G,S}(r)$ for the cardinality of the ball of radius $r$ in the Cayley graph $\Cay(G,S)$.

\begin{thm}[Breuillard, Green, and Tao 2012]
\label{thm:BGT_metric}
For each $K\geq 1$ there exist constants $r_0=r_0(K)$ and $C=C(K)$ such that the following holds. Let $G$ be a group with finite generating set $S$, and suppose that there exists $r \geq r_0$ such that $\Gr(3r) \leq K\Gr(r)$. Then $\Gr(mr) \leq m^C \Gr(r)$ for every $m\geq 3$ and there exists a finite normal subgroup $Q \triangleleft G$ such that:

\begin{enumerate}
    \item Every fibre of the projection $\pi:G \rightarrow G/Q$ has diameter at most $Cr$.
    
    %
    
    \item $G/Q$ has a nilpotent normal subgroup $N$ of rank, step and index at most $C$.
    
    %
    
    %
    
    %
    \item The projection $x\mapsto \pi(x)$ is a $(1,Cr)$-quasi-isometry from $\operatorname{Cay}(G,S)$ to $\operatorname{Cay}(G/Q,\pi(S))$.
\end{enumerate}
\end{thm}

Here, we recall that a function $\phi:V_1\to V_2$ between the vertex sets of two graphs $G_1=(V_1,E_1)$ and $G_2=(V_2,E_2)$ is said to be an $(\alpha,\beta)$\textbf{-quasi-isometry} (a.k.a.\ \textbf{rough isometry}) if 
\[\alpha^{-1} d(x,y) -\beta \leq d(\phi(x),\phi(y)) \leq \alpha d(x,y) +\beta\]
for every $x,y\in V_1$, and every vertex $z\in V_2$ is within distance at most $\beta$ of $\phi(V_1)$. (The second property holds automatically if $\phi$ is surjective.)

\medskip

Informally, the Breuillard-Green-Tao theorem states that polynomial growth at one sufficiently large scale forces the group to have polynomial growth at every subsequent scale, and moreover to be metrically ``well-modelled'' by a nilpotent group at all larger scales. Similar theorems for vertex-transitive graphs that are not necessarily Cayley graphs have recently been established in the work of Tessera and Tointon \cite{tt.Trof,MR3877012}.
